\documentclass[10pt,a4paper]{extarticle}
\usepackage[utf8]{inputenc}
\usepackage[T1]{fontenc}
\usepackage[spanish,es-tabla,es-noshorthands]{babel}
\usepackage[top=1.6cm,bottom=1.4cm,left=1.3cm,right=1.3cm,headheight=11pt,headsep=6pt]{geometry}
\usepackage{amsmath,amssymb,mathtools}
\usepackage{icomma}
\spanishplainpercent
\usepackage{booktabs,array,multicol,enumitem}
\usepackage[table]{xcolor}
\usepackage[most]{tcolorbox}
\usepackage{fancyhdr}

\definecolor{azul}{RGB}{35,85,140}
\definecolor{azulclaro}{RGB}{238,244,251}
\definecolor{rojo}{RGB}{180,40,40}

\setlength{\parindent}{0pt}\setlength{\tabcolsep}{3pt}
\setlength{\parskip}{2pt}
\setlist{itemsep=0pt,topsep=1pt,leftmargin=10pt}
\pagestyle{fancy}\fancyhf{}
\lhead{\footnotesize\textbf{Hoja de fórmulas} --- Probabilidad y Estadística (Unidades 5, 6 y 7)}
\rhead{\footnotesize UTN FRH}
\cfoot{\footnotesize\thepage}
\renewcommand{\headrulewidth}{0.3pt}

\newtcolorbox{hoja}[2][]{colback=white,colframe=azul!70,boxrule=0.5pt,arc=1.5pt,left=3pt,right=3pt,top=1.5pt,bottom=1.5pt,
  fonttitle=\bfseries\footnotesize,colbacktitle=azul!85,coltitle=white,title={#2},boxsep=1pt,#1}

\newcommand{\E}{\mathrm{E}}
\newcommand{\V}{\mathrm{V}}
\newcommand{\Cov}{\mathrm{Cov}}
\newcommand{\Nor}{\mathcal{N}}
\newcommand{\xb}{\bar{x}}
\newcommand{\Xb}{\bar{X}}
\newcommand{\chis}{\chi^2}
\newcommand{\sq}[1]{\sqrt{#1}}

\begin{document}
\footnotesize
\begin{center}
{\Large\bfseries\color{azul} Hoja de fórmulas --- Probabilidad y Estadística}\\[1pt]
{\footnotesize Variables bidimensionales $\cdot$ Estimación (IC) $\cdot$ Pruebas de hipótesis \quad|\quad
$\Phi(z)=P(Z\le z)$; \ $z_\alpha,\,t_{\alpha;\nu},\,\chis_{\alpha;\nu}$ dejan área $\alpha$ \emph{a la derecha}; \ $q=1-p$}
\end{center}
\vspace{-4pt}

% ============================================================ PAGINA 1
\begin{hoja}{A. Distribuciones y propiedades básicas}
\renewcommand{\arraystretch}{1.3}
\begin{tabular}{@{}>{\raggedright\arraybackslash}p{3.7cm}>{\raggedright\arraybackslash}p{6.6cm}>{\raggedright\arraybackslash}p{2.0cm}>{\raggedright\arraybackslash}p{5.0cm}@{}}
\toprule
\textbf{Distribución} & \textbf{$p(x)$ o $f(x)$} & $\E(X)$ & $\V(X)$\\\midrule
Bernoulli$(p)$ & $p^xq^{1-x}$, $x\in\{0,1\}$ & $p$ & $pq$\\
Binomial$(n;p)$ & $\binom nxp^xq^{n-x}$, $x=0..n$ & $np$ & $npq$\\
Hipergeométrica$(N;M;n)$ & $\binom Mx\binom{N-M}{n-x}\big/\binom Nn$ & $n\frac MN$ & $n\frac MN\big(1-\frac MN\big)\frac{N-n}{N-1}$\\
Poisson$(\lambda)$ & $e^{-\lambda}\lambda^x/x!$, $x=0,1,\dots$ & $\lambda$ & $\lambda$\\
Geométrica$(p)$ & $q^{x-1}p$, $x=1,2,\dots$ & $1/p$ & $q/p^2$\\
Uniforme$(a;b)$ & $\frac1{b-a}$, $a<x<b$ & $\frac{a+b}2$ & $\frac{(b-a)^2}{12}$\\
Exponencial$(\lambda)$ & $\lambda e^{-\lambda x}$, $x\ge0$ & $1/\lambda$ & $1/\lambda^2$\\
Normal $\Nor(\mu;\sigma^2)$ & $\frac1{\sigma\sq{2\pi}}e^{-(x-\mu)^2/(2\sigma^2)}$ & $\mu$ & $\sigma^2$\\
$t_\nu$ \ / \ $\chis_\nu$ & $t$: simétrica en 0; \ $\chis$: $\ge0$ y asimétrica & $0$ \ / \ $\nu$ & $\frac{\nu}{\nu-2}$ ($\nu>2$) \ / \ $2\nu$\\
\bottomrule
\end{tabular}

\smallskip
$\E(aX+b)=a\E(X)+b$; \ $\V(aX+b)=a^2\V(X)$; \ $\V(X)=\E(X^2)-\E(X)^2$. \quad
$Z=\frac{X-\mu}{\sigma}\sim\Nor(0;1)$; \ $P(a<X<b)=\Phi\big(\frac{b-\mu}\sigma\big)-\Phi\big(\frac{a-\mu}\sigma\big)$; \ $\Phi(-z)=1-\Phi(z)$; \ percentil: $x=\mu+z\sigma$.\\
Chebyshev (cualquier distribución): $P(|X-\mu|\ge k\sigma)\le\frac1{k^2}$.
\end{hoja}

\begin{hoja}{B. Variables aleatorias bidimensionales}
\begin{multicols}{2}\raggedright
\textbf{Discreto.} $p(x,y)=P(X{=}x,Y{=}y)\ge0$, $\sum\sum p=1$.\\
$p_X(x)=\sum_yp(x,y)$; \ $p_Y(y)=\sum_xp(x,y)$.\\
$p_{Y|X}(y|x)=\frac{p(x,y)}{p_X(x)}$; \ $\E(Y|x)=\sum_yy\,p_{Y|X}(y|x)$.\\
$P(A|B)=\frac{P(A\cap B)}{P(B)}$; \ $P(A\cap B)=P(A)P(B|A)$.\\
Conteo sin reposición: $p(x,y)=\frac{\binom{N_1}x\binom{N_2}y\binom{N_3}{n-x-y}}{\binom Nn}$.\\
\textbf{Independientes} $\iff p(x,y)=p_X(x)p_Y(y)$ $\forall(x,y)$.\\
$\E[h(X,Y)]=\sum\sum h(x,y)p(x,y)$; \ $\E(XY)=\sum\sum xy\,p(x,y)$.

\columnbreak
\textbf{Continuo.} $f\ge0$, $\iint f=1$; \ $P[(X,Y)\in A]=\iint_Af$.\\
$f_X(x)=\int f(x,y)dy$; \ $f_Y(y)=\int f(x,y)dx$ (límites según el soporte).\\
$f_{Y|X}(y|x)=\frac{f(x,y)}{f_X(x)}$; \ $\E[h]=\iint h\,f$.\\
Indep.: $f=f_Xf_Y$ (el soporte debe ser rectángulo).\\
\textbf{Multinomial} $(n;p_1,p_2,\dots)$: \ $\Cov(X_i,X_j)=-np_ip_j$.
\end{multicols}
\vspace{-5pt}
\textbf{Covarianza:} $\Cov(X,Y)=\E(XY)-\mu_X\mu_Y$; \ $\Cov(X,X)=\V(X)$; \ $\Cov(aX+b,cY+d)=ac\Cov(X,Y)$; \ $\Cov(h-X,X)=-\V(X)$.
\ \textbf{Correlación:} $\rho=\frac{\Cov(X,Y)}{\sigma_X\sigma_Y}\in[-1;1]$; $|\rho|=1\iff Y=aX+b$; $\Cov=\rho\sigma_X\sigma_Y$. Independencia $\Rightarrow\Cov=0$ (no al revés).
\end{hoja}

\begin{hoja}{C. Combinaciones lineales, propiedad reproductiva, TCL y aproximaciones}
\raggedright
$W=\sum a_iX_i+b$: \ $\E(W)=\sum a_i\mu_i+b$; \ $\V(W)=\sum a_i^2\sigma_i^2+2\sum_{i<j}a_ia_j\Cov(X_i,X_j)$. \quad
$\V(aX+bY+c)=a^2\V(X)+b^2\V(Y)+2ab\Cov(X,Y)$; \ independientes: $\V(X\pm Y)=\V(X)+\V(Y)$.\\
iid: $T=\sum X_i$: $\E=n\mu$, $\V=n\sigma^2$; \ $\Xb$: $\E=\mu$, $\V=\sigma^2/n$ (EE $=\sigma/\sq n$). \ Población finita: $\V(\Xb)=\frac{\sigma^2}n\frac{N-n}{N-1}$.\\
\textbf{Reproductiva} (exacta): $X_i$ normales indep.\ $\Rightarrow\sum a_iX_i+b\sim\Nor\big(\sum a_i\mu_i+b;\ \sum a_i^2\sigma_i^2\big)$. \
\textbf{TCL} (aprox., $n\ge30$, cualquier población): $\Xb\approx\Nor(\mu;\frac{\sigma^2}n)$, \ $\sum X_i\approx\Nor(n\mu;n\sigma^2)$.\\
\textbf{Binomial $\to$ normal} ($np\ge5$, $nq\ge5$): $X\approx\Nor(np;npq)$, $\hat P\approx\Nor(p;\frac{pq}n)$. \textbf{Corrección por continuidad:} $P(X\le k)\to k{+}0{,}5$; \ $P(X<k)\to k{-}0{,}5$; \ $P(X\ge k)\to k{-}0{,}5$; \ $P(X>k)\to k{+}0{,}5$; \ $P(X=k)\to[k{-}0{,}5;k{+}0{,}5]$.\\
\textbf{Binomial $\to$ Poisson} (evento raro): $\lambda=np$. \quad Sumas indep.: $\sum\mathrm{Bin}(n_i;p)=\mathrm{Bin}(\sum n_i;p)$; \ $\sum\mathrm{Poi}(\lambda_i)=\mathrm{Poi}(\sum\lambda_i)$; \ $\sum\chis_{\nu_i}=\chis_{\sum\nu_i}$.
\end{hoja}

\begin{hoja}{D. Muestreo y distribuciones muestrales}
\raggedright
$\xb=\frac1n\sum x_i$; \ $s^2=\frac{\sum(x_i-\xb)^2}{n-1}=\frac{\sum x_i^2-n\xb^2}{n-1}$; \ $\hat p=\frac xn$; \ $\sum(x_i-\xb)^2=(n-1)s^2$.\\[1pt]
\renewcommand{\arraystretch}{1.55}
\begin{tabular}{@{}>{\raggedright\arraybackslash}p{5.2cm}>{\raggedright\arraybackslash}p{3.9cm}>{\raggedright\arraybackslash}p{8.2cm}@{}}
\toprule
\textbf{Estadístico} & \textbf{Distribución} & \textbf{Condición}\\\midrule
$Z=\dfrac{\Xb-\mu}{\sigma/\sq n}$ & $\Nor(0;1)$ & población normal (exacta) o $n\ge30$ (TCL)\\
$T=\dfrac{\Xb-\mu}{S/\sq n}$ & $t_{n-1}$ & población normal, $\sigma$ desconocido\\
$\chis=\dfrac{(n-1)S^2}{\sigma^2}$ & $\chis_{n-1}$ & población normal (no es robusto)\\
$Z=\dfrac{\hat P-p}{\sq{pq/n}}$ & $\approx\Nor(0;1)$ & $np\ge5$ y $nq\ge5$\\
$\Xb_1-\Xb_2$ & $\Nor\big(\mu_1-\mu_2;\frac{\sigma_1^2}{n_1}+\frac{\sigma_2^2}{n_2}\big)$ & muestras independientes, poblaciones normales (o $n_i\ge30$)\\
$T=\dfrac{(\Xb_1-\Xb_2)-(\mu_1-\mu_2)}{S_p\sq{1/n_1+1/n_2}}$ & $t_{n_1+n_2-2}$ & normales, $\sigma_1=\sigma_2$ desconocidas; \ $S_p^2=\frac{(n_1-1)S_1^2+(n_2-1)S_2^2}{n_1+n_2-2}$\\
$T=\dfrac{\bar D-\mu_D}{S_D/\sq n}$ & $t_{n-1}$ & datos apareados ($D_i=X_{1i}-X_{2i}$ normales; $n$ = pares)\\
\bottomrule
\end{tabular}
\end{hoja}

\newpage
% ============================================================ PAGINA 2
\begin{hoja}{E. Estimación puntual}
\raggedright
Sesgo $b=\E(\hat\Theta)-\theta$ (insesgado si $b=0$). \ $\mathrm{ECM}=\E[(\hat\Theta-\theta)^2]=\V(\hat\Theta)+b^2$. \ Eficiencia relativa $=\V(\hat\Theta_1)/\V(\hat\Theta_2)$; entre insesgados, mejor el de menor varianza. \ Consistente: $\mathrm{ECM}\to0$ si $n\to\infty$.\\
Insesgados: $\Xb\to\mu$ (EE $\sigma/\sq n$); \ $\hat P\to p$ (EE $\sq{pq/n}$); \ $S^2\to\sigma^2$ ($S$ es sesgado para $\sigma$); \ $\Xb_1-\Xb_2$; \ $\hat P_1-\hat P_2$; \ $S_p^2$. \ Combinación óptima de dos medias: $a\Xb_1+(1-a)\Xb_2$ con $a=\frac{n_1}{n_1+n_2}$. \ $\E(X^2)=\V(X)+\E(X)^2$.
\end{hoja}

\begin{hoja}{F. Intervalos de confianza del $(1-\alpha)100\,\%$ \quad ($E$ = error máximo = semiamplitud = valor crítico $\times$ error estándar)}
\renewcommand{\arraystretch}{1.5}
\begin{tabular}{@{}>{\raggedright\arraybackslash}p{1.7cm}>{\raggedright\arraybackslash}p{4.6cm}>{\raggedright\arraybackslash}p{10.9cm}@{}}
\toprule
\textbf{Parámetro} & \textbf{Condiciones} & \textbf{Intervalo}\\\midrule
$\mu$ & $\sigma$ conocido (normal o $n\ge30$) & $\xb\pm z_{\alpha/2}\dfrac{\sigma}{\sq n}$\\
$\mu$ & $\sigma$ desconocido, normal & $\xb\pm t_{\alpha/2;\,n-1}\dfrac{s}{\sq n}$\\
$\mu$ & $\sigma$ desconocido, $n\ge30$ & $\xb\pm z_{\alpha/2}\dfrac{s}{\sq n}$\\
$\sigma^2$ \ ($\sigma$) & población normal & $\Big[\dfrac{(n-1)s^2}{\chis_{\alpha/2;\,n-1}}\ ;\ \dfrac{(n-1)s^2}{\chis_{1-\alpha/2;\,n-1}}\Big]$ \quad (para $\sigma$: raíz de los extremos; el valor crítico \emph{grande} va abajo)\\
$p$ & $n\hat p\ge5$, $n\hat q\ge5$ & $\hat p\pm z_{\alpha/2}\sq{\dfrac{\hat p\hat q}{n}}$\\
$\mu_1-\mu_2$ & $\sigma_i$ conocidos, o $n_i\ge30$ (con $s_i$) & $(\xb_1-\xb_2)\pm z_{\alpha/2}\sq{\dfrac{\sigma_1^2}{n_1}+\dfrac{\sigma_2^2}{n_2}}$\\
$\mu_1-\mu_2$ & normales, $\sigma_1=\sigma_2$ desconocidos & $(\xb_1-\xb_2)\pm t_{\alpha/2;\,n_1+n_2-2}\;s_p\sq{\dfrac1{n_1}+\dfrac1{n_2}}$, \quad $s_p^2=\dfrac{(n_1-1)s_1^2+(n_2-1)s_2^2}{n_1+n_2-2}$\\
$\mu_1-\mu_2$ & normales, $\sigma_1\ne\sigma_2$ (Welch) & $(\xb_1-\xb_2)\pm t_{\alpha/2;\,\nu}\sq{\dfrac{s_1^2}{n_1}+\dfrac{s_2^2}{n_2}}$, \quad $\nu=\dfrac{(s_1^2/n_1+s_2^2/n_2)^2}{\frac{(s_1^2/n_1)^2}{n_1-1}+\frac{(s_2^2/n_2)^2}{n_2-1}}$ (redondear hacia abajo)\\
$\mu_D$ & datos apareados & $\bar d\pm t_{\alpha/2;\,n-1}\dfrac{s_d}{\sq n}$ \quad ($d_i=x_{1i}-x_{2i}$)\\
$p_1-p_2$ & muestras grandes & $(\hat p_1-\hat p_2)\pm z_{\alpha/2}\sq{\dfrac{\hat p_1\hat q_1}{n_1}+\dfrac{\hat p_2\hat q_2}{n_2}}$\\
\bottomrule
\end{tabular}

\smallskip
\textbf{Cotas unilaterales} (usar $\alpha$ en lugar de $\alpha/2$): \ $\mu\le\xb+z_\alpha\frac\sigma{\sq n}$; \ $\mu\ge\xb-z_\alpha\frac\sigma{\sq n}$ (igual con $t$); \ $\sigma^2\le\frac{(n-1)s^2}{\chis_{1-\alpha;n-1}}$; \ $\sigma^2\ge\frac{(n-1)s^2}{\chis_{\alpha;n-1}}$; \ $p\le\hat p+z_\alpha\sq{\hat p\hat q/n}$.\\
\textbf{Lectura:} si el IC de una diferencia contiene al 0, no hay evidencia de diferencia. \textbf{Nivel de un IC dado:} $\hat p$ = punto medio, $E$ = semiamplitud, $z=E\big/\sq{\hat p\hat q/n}$, \ $1-\alpha=2\Phi(z)-1$ (para $\mu$: $z=E\sq n/\sigma$).\\
\textbf{Interpretación:} ``con $(1-\alpha)100\,\%$ de confianza el parámetro está entre\ldots''; el $(1-\alpha)$ es la proporción de intervalos que aciertan a la larga, no la probabilidad del intervalo calculado.
\end{hoja}

\begin{hoja}{G. Tamaño de muestra (redondear siempre hacia arriba)}
\renewcommand{\arraystretch}{1.6}
\begin{tabular}{@{}>{\raggedright\arraybackslash}p{3.6cm}>{\raggedright\arraybackslash}p{14.2cm}@{}}
\textbf{Estimar $\mu$} & $n=\Big(\dfrac{z_{\alpha/2}\,\sigma}{E}\Big)^2$ \quad ($E$ = semiamplitud; si dan la amplitud $A$: $E=A/2$; $\sigma$ por muestra piloto o dato histórico)\\
\textbf{Estimar $p$} & $n=\dfrac{z_{\alpha/2}^2\,\hat p\hat q}{E^2}$ (con estimación previa) \quad o \quad $n=\dfrac{z_{\alpha/2}^2}{4E^2}$ (sin dato, o para garantizarlo para cualquier $p$; $p(1-p)\le\frac14$)\\
\textbf{Cambios relativos} & $E\propto z/\sq n$: \ $\sqrt{n_2}=\sqrt{n_1}\,\dfrac{z_2}{z_1}\dfrac{E_1}{E_2}$ ($\sigma$ se cancela). Error a la mitad: $n\times4$. \ ``¿Cuántos más?'': $n_2-n_1$.\\
\textbf{Prueba sobre $\mu$} \ ($\alpha$ y $\beta$ dados) & $n=\Big[\dfrac{(z_\alpha+z_\beta)\,\sigma}{|\mu_0-\mu_1|}\Big]^2$ (unilateral); \ bilateral: $z_{\alpha/2}$ en lugar de $z_\alpha$ (aprox.). \ Solo $\alpha$ y una regla fija $c$: \ $\dfrac{|c-\mu_0|}{\sigma/\sq n}=z_\alpha$.\\
\textbf{Prueba sobre $p$} & $n=\Big[\dfrac{z_\alpha\sq{p_0q_0}+z_\beta\sq{p_1q_1}}{|p_0-p_1|}\Big]^2$ \quad (bilateral: $z_{\alpha/2}$)\\
\end{tabular}
\end{hoja}

\newpage
% ============================================================ PAGINA 3
\begin{hoja}{H. Pruebas de hipótesis}
\raggedright
\textbf{Errores:} \ tipo I $=$ rechazar $H_0$ siendo verdadera, $\alpha=P(\text{rechazar }H_0|H_0\text{ V})$; \ tipo II $=$ \textbf{no} rechazar $H_0$ siendo falsa, $\beta=P(\text{no rechazar }H_0|H_0\text{ F})$; \ potencia $=1-\beta=P(\text{rechazar }H_0|H_0\text{ F})$. Con $n$ fijo, $\downarrow\alpha\Rightarrow\uparrow\beta$; al aumentar $n$ bajan los dos. Si se rechazó $H_0$ el error posible es el I; si no se rechazó, el II.\\[2pt]
\renewcommand{\arraystretch}{1.6}
\begin{tabular}{@{}>{\raggedright\arraybackslash}p{2.5cm}>{\raggedright\arraybackslash}p{4.7cm}>{\raggedright\arraybackslash}p{6.6cm}>{\raggedright\arraybackslash}p{3.0cm}@{}}
\toprule
\textbf{$H_0$} & \textbf{Condiciones} & \textbf{Estadístico de prueba} & \textbf{Distribución}\\\midrule
$\mu=\mu_0$ & $\sigma$ conocido (normal o $n\ge30$) & $Z_0=\dfrac{\xb-\mu_0}{\sigma/\sq n}$ & $\Nor(0;1)$\\
$\mu=\mu_0$ & normal, $\sigma$ desconocido & $T_0=\dfrac{\xb-\mu_0}{s/\sq n}$ & $t_{n-1}$\\
$\sigma^2=\sigma_0^2$ & población normal & $\chis_0=\dfrac{(n-1)s^2}{\sigma_0^2}$ & $\chis_{n-1}$\\
$p=p_0$ & $np_0\ge5$, $nq_0\ge5$ & $Z_0=\dfrac{\hat p-p_0}{\sq{p_0q_0/n}}$ \ (¡con $p_0$!) & $\Nor(0;1)$\\
$p=p_0$ & $n$ chico & $X$ = cantidad de éxitos; $\alpha=P(X\in R\,|\,p_0)$ & $\mathrm{Bin}(n;p_0)$\\
$\mu_1-\mu_2=\Delta_0$ & $\sigma_i$ conocidos o $n_i\ge30$ & $Z_0=\dfrac{(\xb_1-\xb_2)-\Delta_0}{\sq{\sigma_1^2/n_1+\sigma_2^2/n_2}}$ & $\Nor(0;1)$\\
$\mu_1-\mu_2=\Delta_0$ & normales, $\sigma_1=\sigma_2$ & $T_0=\dfrac{(\xb_1-\xb_2)-\Delta_0}{s_p\sq{1/n_1+1/n_2}}$ & $t_{n_1+n_2-2}$\\
$\mu_1-\mu_2=\Delta_0$ & normales, $\sigma_1\ne\sigma_2$ & $T_0=\dfrac{(\xb_1-\xb_2)-\Delta_0}{\sq{s_1^2/n_1+s_2^2/n_2}}$ & $t_\nu$ (Welch)\\
$\mu_D=\Delta_0$ & apareadas & $T_0=\dfrac{\bar d-\Delta_0}{s_d/\sq n}$ & $t_{n-1}$ ($n$ = pares)\\
$p_1=p_2$ & muestras grandes & $Z_0=\dfrac{\hat p_1-\hat p_2}{\sq{\hat p\hat q(1/n_1+1/n_2)}}$, \ $\hat p=\frac{x_1+x_2}{n_1+n_2}$ & $\Nor(0;1)$\\
\bottomrule
\end{tabular}

\smallskip
\renewcommand{\arraystretch}{1.5}
\begin{tabular}{@{}>{\raggedright\arraybackslash}p{1.4cm}>{\raggedright\arraybackslash}p{2.9cm}>{\raggedright\arraybackslash}p{4.6cm}>{\raggedright\arraybackslash}p{2.9cm}>{\raggedright\arraybackslash}p{5.2cm}@{}}
\toprule
$H_1$ & Rechazo ($Z$; $t$ igual con $\nu$) & Rechazo ($\chis$, $\nu=n-1$) & Valor $P$ ($Z$) & $\beta(\mu_1)$ \ con \ $d=\dfrac{(\mu_1-\mu_0)\sq n}{\sigma}$\\\midrule
$\theta\neq\theta_0$ & $|z_0|>z_{\alpha/2}$ & $\chis_0>\chis_{\alpha/2}$ \ o \ $\chis_0<\chis_{1-\alpha/2}$ & $2[1-\Phi(|z_0|)]$ & $\Phi(z_{\alpha/2}-d)-\Phi(-z_{\alpha/2}-d)$\\
$\theta>\theta_0$ & $z_0>z_\alpha$ & $\chis_0>\chis_{\alpha}$ & $1-\Phi(z_0)$ & $\Phi(z_\alpha-d)$\\
$\theta<\theta_0$ & $z_0<-z_\alpha$ & $\chis_0<\chis_{1-\alpha}$ & $\Phi(z_0)$ & $1-\Phi(-z_\alpha-d)$\\
\bottomrule
\end{tabular}

\smallskip
\textbf{Decisión:} $P\le\alpha\Rightarrow$ se rechaza $H_0$ (el valor $P$ no es la probabilidad de que $H_0$ sea cierta). Con tablas $t$ o $\chis$ se acota $P$ entre dos columnas (bilateral: $\times2$; $\chis$ bilateral: $P=2\min\{P(\chis\ge\chis_0),P(\chis\le\chis_0)\}$).\\
\textbf{Región en unidades originales:} $c=\mu_0\pm z\frac{\sigma}{\sq n}$ \ ($p$: $c=p_0\pm z\sq{p_0q_0/n}$). \ \textbf{$\alpha$ de una regla dada:} $\alpha=P(\Xb\in R\,|\,\mu_0)$.\\
\textbf{$\beta$ en dos pasos:} (1) región de aceptación en $\xb$ (o $\hat p$) calculada bajo $H_0$; (2) $\beta=P(\Xb\in\text{aceptación}\,|\,\mu_1)$, con la normal centrada en $\mu_1$ y el mismo $\sigma/\sq n$. \ \textbf{Para $p$:} desvío $\sq{p_1q_1/n}$ y, p.ej.\ $H_1{:}\,p<p_0$: \ $\beta=1-\Phi\Big(\dfrac{p_0-p_1-z_\alpha\sq{p_0q_0/n}}{\sq{p_1q_1/n}}\Big)$; \ $H_1{:}\,p>p_0$: \ $\beta=\Phi\Big(\dfrac{p_0-p_1+z_\alpha\sq{p_0q_0/n}}{\sq{p_1q_1/n}}\Big)$.\\
$\delta=\mu_1-\mu_0$ es la diferencia a detectar (con $H_1{:}\,\mu<\mu_0$, $\delta<0$ y $\mu_1$ del lado de $H_1$). \ Potencia $\pi(\mu_1)=1-\beta(\mu_1)$.\\
\textbf{Curvas:} la OC $\beta(\theta)$ vale $1-\alpha$ en $\theta_0$ y la potencia vale $\alpha$. Bilateral: OC en campana (máx.\ $1-\alpha$), potencia en valle. $H_1{:}>$: OC decreciente, potencia creciente. $H_1{:}<$: al revés.\\
\textbf{Dualidad IC--prueba:} una prueba bilateral con nivel $\alpha$ rechaza $H_0{:}\,\theta=\theta_0$ $\iff$ $\theta_0\notin$ IC$(1-\alpha)$. La unilateral se compara con la cota del mismo lado. Para proporciones es solo aproximado ($p_0$ en la prueba, $\hat p$ en el IC).
\end{hoja}

\begin{hoja}{I. Cómo plantear $H_0$ y $H_1$, y los 7 pasos de toda prueba}
\begin{multicols}{2}\raggedright
\textbf{$H_0$ siempre con ``$=$''} (el statu quo, lo que se afirma); \textbf{$H_1$} = lo que se quiere demostrar o lo que lleva a actuar.\\
``supera'', ``aumentó'', ``excede'', ``a lo sumo (¿hay evidencia en contra?)'' $\Rightarrow$ $H_1{:}\,>$.\\
``es menor'', ``disminuyó'', ``afirma que \emph{al menos}\ldots (actuará si se prueba que no)'' $\Rightarrow$ $H_1{:}\,<$.\\
``difiere'', ``cambió'', ``¿hay que ajustar?'' $\Rightarrow$ $H_1{:}\,\ne$.\\
``se aprueba / firma / incorpora solo si se prueba que cumple'' $\Rightarrow$ $H_1$ = cumple.

\columnbreak
\textbf{(1)} parámetro y datos \ \textbf{(2)} $H_0$ \ \textbf{(3)} $H_1$ \ \textbf{(4)} $\alpha$ (0,05 si no lo dan) \ \textbf{(5)} estadístico y su distribución bajo $H_0$, justificando (normalidad, $\sigma$, TCL) \ \textbf{(6)} región de rechazo \ \textbf{(7)} calcular, decidir y concluir en el contexto: ``se rechaza $H_0$: hay evidencia de que\ldots'' o ``no se rechaza: no hay evidencia de que\ldots'' (nunca ``se acepta $H_0$'').\\
\textbf{Controles:} las tablas suman 1; $-1\le\rho\le1$; $P\in[0;1]$; en $\theta_0$ la OC vale $1-\alpha$; $n$ siempre hacia arriba; usar al menos 4 decimales.
\end{multicols}
\end{hoja}

\begin{hoja}{J. Valores críticos más usados}
\renewcommand{\arraystretch}{1.25}
\begin{tabular}{@{}lcccccccccc@{}}
\toprule
$\alpha$ & 0,20 & 0,10 & 0,06 & 0,05 & 0,025 & 0,02 & 0,015 & 0,01 & 0,005 & 0,001\\\midrule
$z_\alpha$ & 0,8416 & 1,2816 & 1,5548 & 1,6449 & 1,9600 & 2,0537 & 2,1701 & 2,3263 & 2,5758 & 3,0902\\\bottomrule
\end{tabular}\\[2pt]
IC bilateral: 90\,\% $\to1{,}6449$; \ 95\,\% $\to1{,}96$; \ 98\,\% $\to2{,}3263$; \ 99\,\% $\to2{,}5758$. \quad $\Phi(1)=0{,}8413$; \ $\Phi(2)=0{,}9772$; \ $\Phi(3)=0{,}9987$.
\end{hoja}

% ============================================================ ANEXO
\newpage
\begin{center}{\large\bfseries\color{azul} Anexo --- Tablas}\end{center}
\vspace{-6pt}
\begin{multicols}{2}
\textbf{Normal estándar: $\Phi(z)=P(Z\le z)$.} Para $z<0$: $\Phi(-z)=1-\Phi(z)$.
\begin{center}
\scriptsize\setlength{\tabcolsep}{1.5pt}\renewcommand{\arraystretch}{0.92}
\begin{tabular}{@{}c|*{10}{c}@{}}\toprule
$z$ & 0,00 & 0,01 & 0,02 & 0,03 & 0,04 & 0,05 & 0,06 & 0,07 & 0,08 & 0,09\\\midrule
0,0 & 0,5000 & 0,5040 & 0,5080 & 0,5120 & 0,5160 & 0,5199 & 0,5239 & 0,5279 & 0,5319 & 0,5359\\
0,1 & 0,5398 & 0,5438 & 0,5478 & 0,5517 & 0,5557 & 0,5596 & 0,5636 & 0,5675 & 0,5714 & 0,5753\\
0,2 & 0,5793 & 0,5832 & 0,5871 & 0,5910 & 0,5948 & 0,5987 & 0,6026 & 0,6064 & 0,6103 & 0,6141\\
0,3 & 0,6179 & 0,6217 & 0,6255 & 0,6293 & 0,6331 & 0,6368 & 0,6406 & 0,6443 & 0,6480 & 0,6517\\
0,4 & 0,6554 & 0,6591 & 0,6628 & 0,6664 & 0,6700 & 0,6736 & 0,6772 & 0,6808 & 0,6844 & 0,6879\\
\addlinespace[2pt]
0,5 & 0,6915 & 0,6950 & 0,6985 & 0,7019 & 0,7054 & 0,7088 & 0,7123 & 0,7157 & 0,7190 & 0,7224\\
0,6 & 0,7257 & 0,7291 & 0,7324 & 0,7357 & 0,7389 & 0,7422 & 0,7454 & 0,7486 & 0,7517 & 0,7549\\
0,7 & 0,7580 & 0,7611 & 0,7642 & 0,7673 & 0,7704 & 0,7734 & 0,7764 & 0,7794 & 0,7823 & 0,7852\\
0,8 & 0,7881 & 0,7910 & 0,7939 & 0,7967 & 0,7995 & 0,8023 & 0,8051 & 0,8078 & 0,8106 & 0,8133\\
0,9 & 0,8159 & 0,8186 & 0,8212 & 0,8238 & 0,8264 & 0,8289 & 0,8315 & 0,8340 & 0,8365 & 0,8389\\
\addlinespace[2pt]
1,0 & 0,8413 & 0,8438 & 0,8461 & 0,8485 & 0,8508 & 0,8531 & 0,8554 & 0,8577 & 0,8599 & 0,8621\\
1,1 & 0,8643 & 0,8665 & 0,8686 & 0,8708 & 0,8729 & 0,8749 & 0,8770 & 0,8790 & 0,8810 & 0,8830\\
1,2 & 0,8849 & 0,8869 & 0,8888 & 0,8907 & 0,8925 & 0,8944 & 0,8962 & 0,8980 & 0,8997 & 0,9015\\
1,3 & 0,9032 & 0,9049 & 0,9066 & 0,9082 & 0,9099 & 0,9115 & 0,9131 & 0,9147 & 0,9162 & 0,9177\\
1,4 & 0,9192 & 0,9207 & 0,9222 & 0,9236 & 0,9251 & 0,9265 & 0,9279 & 0,9292 & 0,9306 & 0,9319\\
\addlinespace[2pt]
1,5 & 0,9332 & 0,9345 & 0,9357 & 0,9370 & 0,9382 & 0,9394 & 0,9406 & 0,9418 & 0,9429 & 0,9441\\
1,6 & 0,9452 & 0,9463 & 0,9474 & 0,9484 & 0,9495 & 0,9505 & 0,9515 & 0,9525 & 0,9535 & 0,9545\\
1,7 & 0,9554 & 0,9564 & 0,9573 & 0,9582 & 0,9591 & 0,9599 & 0,9608 & 0,9616 & 0,9625 & 0,9633\\
1,8 & 0,9641 & 0,9649 & 0,9656 & 0,9664 & 0,9671 & 0,9678 & 0,9686 & 0,9693 & 0,9699 & 0,9706\\
1,9 & 0,9713 & 0,9719 & 0,9726 & 0,9732 & 0,9738 & 0,9744 & 0,9750 & 0,9756 & 0,9761 & 0,9767\\
\addlinespace[2pt]
2,0 & 0,9772 & 0,9778 & 0,9783 & 0,9788 & 0,9793 & 0,9798 & 0,9803 & 0,9808 & 0,9812 & 0,9817\\
2,1 & 0,9821 & 0,9826 & 0,9830 & 0,9834 & 0,9838 & 0,9842 & 0,9846 & 0,9850 & 0,9854 & 0,9857\\
2,2 & 0,9861 & 0,9864 & 0,9868 & 0,9871 & 0,9875 & 0,9878 & 0,9881 & 0,9884 & 0,9887 & 0,9890\\
2,3 & 0,9893 & 0,9896 & 0,9898 & 0,9901 & 0,9904 & 0,9906 & 0,9909 & 0,9911 & 0,9913 & 0,9916\\
2,4 & 0,9918 & 0,9920 & 0,9922 & 0,9925 & 0,9927 & 0,9929 & 0,9931 & 0,9932 & 0,9934 & 0,9936\\
\addlinespace[2pt]
2,5 & 0,9938 & 0,9940 & 0,9941 & 0,9943 & 0,9945 & 0,9946 & 0,9948 & 0,9949 & 0,9951 & 0,9952\\
2,6 & 0,9953 & 0,9955 & 0,9956 & 0,9957 & 0,9959 & 0,9960 & 0,9961 & 0,9962 & 0,9963 & 0,9964\\
2,7 & 0,9965 & 0,9966 & 0,9967 & 0,9968 & 0,9969 & 0,9970 & 0,9971 & 0,9972 & 0,9973 & 0,9974\\
2,8 & 0,9974 & 0,9975 & 0,9976 & 0,9977 & 0,9977 & 0,9978 & 0,9979 & 0,9979 & 0,9980 & 0,9981\\
2,9 & 0,9981 & 0,9982 & 0,9982 & 0,9983 & 0,9984 & 0,9984 & 0,9985 & 0,9985 & 0,9986 & 0,9986\\
\addlinespace[2pt]
3,0 & 0,9987 & 0,9987 & 0,9987 & 0,9988 & 0,9988 & 0,9989 & 0,9989 & 0,9989 & 0,9990 & 0,9990\\
3,1 & 0,9990 & 0,9991 & 0,9991 & 0,9991 & 0,9992 & 0,9992 & 0,9992 & 0,9992 & 0,9993 & 0,9993\\
3,2 & 0,9993 & 0,9993 & 0,9994 & 0,9994 & 0,9994 & 0,9994 & 0,9994 & 0,9995 & 0,9995 & 0,9995\\
3,3 & 0,9995 & 0,9995 & 0,9995 & 0,9996 & 0,9996 & 0,9996 & 0,9996 & 0,9996 & 0,9996 & 0,9997\\
3,4 & 0,9997 & 0,9997 & 0,9997 & 0,9997 & 0,9997 & 0,9997 & 0,9997 & 0,9997 & 0,9997 & 0,9998\\
\bottomrule\end{tabular}
\end{center}

\columnbreak
\textbf{$t$ de Student:} $t_{\alpha;\nu}$ con $P(T_\nu>t_{\alpha;\nu})=\alpha$.
\begin{center}
\footnotesize\setlength{\tabcolsep}{4pt}\renewcommand{\arraystretch}{0.95}
\begin{tabular}{@{}c|*{5}{c}@{}}\toprule
$\nu$ & $0,10$ & $0,05$ & $0,025$ & $0,01$ & $0,005$\\\midrule
1 & 3,078 & 6,314 & 12,706 & 31,821 & 63,657\\
2 & 1,886 & 2,920 & 4,303 & 6,965 & 9,925\\
3 & 1,638 & 2,353 & 3,182 & 4,541 & 5,841\\
4 & 1,533 & 2,132 & 2,776 & 3,747 & 4,604\\
5 & 1,476 & 2,015 & 2,571 & 3,365 & 4,032\\
\addlinespace[2pt]
6 & 1,440 & 1,943 & 2,447 & 3,143 & 3,707\\
7 & 1,415 & 1,895 & 2,365 & 2,998 & 3,499\\
8 & 1,397 & 1,860 & 2,306 & 2,896 & 3,355\\
9 & 1,383 & 1,833 & 2,262 & 2,821 & 3,250\\
10 & 1,372 & 1,812 & 2,228 & 2,764 & 3,169\\
\addlinespace[2pt]
11 & 1,363 & 1,796 & 2,201 & 2,718 & 3,106\\
12 & 1,356 & 1,782 & 2,179 & 2,681 & 3,055\\
13 & 1,350 & 1,771 & 2,160 & 2,650 & 3,012\\
14 & 1,345 & 1,761 & 2,145 & 2,624 & 2,977\\
15 & 1,341 & 1,753 & 2,131 & 2,602 & 2,947\\
\addlinespace[2pt]
16 & 1,337 & 1,746 & 2,120 & 2,583 & 2,921\\
17 & 1,333 & 1,740 & 2,110 & 2,567 & 2,898\\
18 & 1,330 & 1,734 & 2,101 & 2,552 & 2,878\\
19 & 1,328 & 1,729 & 2,093 & 2,539 & 2,861\\
20 & 1,325 & 1,725 & 2,086 & 2,528 & 2,845\\
\addlinespace[2pt]
21 & 1,323 & 1,721 & 2,080 & 2,518 & 2,831\\
22 & 1,321 & 1,717 & 2,074 & 2,508 & 2,819\\
23 & 1,319 & 1,714 & 2,069 & 2,500 & 2,807\\
24 & 1,318 & 1,711 & 2,064 & 2,492 & 2,797\\
25 & 1,316 & 1,708 & 2,060 & 2,485 & 2,787\\
\addlinespace[2pt]
26 & 1,315 & 1,706 & 2,056 & 2,479 & 2,779\\
27 & 1,314 & 1,703 & 2,052 & 2,473 & 2,771\\
28 & 1,313 & 1,701 & 2,048 & 2,467 & 2,763\\
29 & 1,311 & 1,699 & 2,045 & 2,462 & 2,756\\
30 & 1,310 & 1,697 & 2,042 & 2,457 & 2,750\\
\addlinespace[2pt]
40 & 1,303 & 1,684 & 2,021 & 2,423 & 2,704\\
60 & 1,296 & 1,671 & 2,000 & 2,390 & 2,660\\
120 & 1,289 & 1,658 & 1,980 & 2,358 & 2,617\\
$\infty$ & 1,282 & 1,645 & 1,960 & 2,326 & 2,576\\
\bottomrule\end{tabular}
\end{center}
\end{multicols}
\vspace{-4pt}
\textbf{Ji-cuadrado:} $\chis_{\alpha;\nu}$ con $P(\chis_\nu>\chis_{\alpha;\nu})=\alpha$ (las columnas de la izquierda, con $\alpha$ grande, dan los valores chicos).
\begin{center}
\footnotesize\setlength{\tabcolsep}{3.6pt}\renewcommand{\arraystretch}{0.95}
\begin{tabular}{@{}c|*{10}{c}@{}}\toprule
$\nu$ & $0,995$ & $0,99$ & $0,975$ & $0,95$ & $0,90$ & $0,10$ & $0,05$ & $0,025$ & $0,01$ & $0,005$\\\midrule
1 & 0,000 & 0,000 & 0,001 & 0,004 & 0,016 & 2,706 & 3,841 & 5,024 & 6,635 & 7,879\\
2 & 0,010 & 0,020 & 0,051 & 0,103 & 0,211 & 4,605 & 5,991 & 7,378 & 9,210 & 10,60\\
3 & 0,072 & 0,115 & 0,216 & 0,352 & 0,584 & 6,251 & 7,815 & 9,348 & 11,34 & 12,84\\
4 & 0,207 & 0,297 & 0,484 & 0,711 & 1,064 & 7,779 & 9,488 & 11,14 & 13,28 & 14,86\\
5 & 0,412 & 0,554 & 0,831 & 1,145 & 1,610 & 9,236 & 11,07 & 12,83 & 15,09 & 16,75\\
\addlinespace[2pt]
6 & 0,676 & 0,872 & 1,237 & 1,635 & 2,204 & 10,64 & 12,59 & 14,45 & 16,81 & 18,55\\
7 & 0,989 & 1,239 & 1,690 & 2,167 & 2,833 & 12,02 & 14,07 & 16,01 & 18,48 & 20,28\\
8 & 1,344 & 1,646 & 2,180 & 2,733 & 3,490 & 13,36 & 15,51 & 17,53 & 20,09 & 21,95\\
9 & 1,735 & 2,088 & 2,700 & 3,325 & 4,168 & 14,68 & 16,92 & 19,02 & 21,67 & 23,59\\
10 & 2,156 & 2,558 & 3,247 & 3,940 & 4,865 & 15,99 & 18,31 & 20,48 & 23,21 & 25,19\\
\addlinespace[2pt]
11 & 2,603 & 3,053 & 3,816 & 4,575 & 5,578 & 17,28 & 19,68 & 21,92 & 24,72 & 26,76\\
12 & 3,074 & 3,571 & 4,404 & 5,226 & 6,304 & 18,55 & 21,03 & 23,34 & 26,22 & 28,30\\
13 & 3,565 & 4,107 & 5,009 & 5,892 & 7,042 & 19,81 & 22,36 & 24,74 & 27,69 & 29,82\\
14 & 4,075 & 4,660 & 5,629 & 6,571 & 7,790 & 21,06 & 23,68 & 26,12 & 29,14 & 31,32\\
15 & 4,601 & 5,229 & 6,262 & 7,261 & 8,547 & 22,31 & 25,00 & 27,49 & 30,58 & 32,80\\
\addlinespace[2pt]
16 & 5,142 & 5,812 & 6,908 & 7,962 & 9,312 & 23,54 & 26,30 & 28,85 & 32,00 & 34,27\\
17 & 5,697 & 6,408 & 7,564 & 8,672 & 10,09 & 24,77 & 27,59 & 30,19 & 33,41 & 35,72\\
18 & 6,265 & 7,015 & 8,231 & 9,390 & 10,86 & 25,99 & 28,87 & 31,53 & 34,81 & 37,16\\
19 & 6,844 & 7,633 & 8,907 & 10,12 & 11,65 & 27,20 & 30,14 & 32,85 & 36,19 & 38,58\\
20 & 7,434 & 8,260 & 9,591 & 10,85 & 12,44 & 28,41 & 31,41 & 34,17 & 37,57 & 40,00\\
\addlinespace[2pt]
21 & 8,034 & 8,897 & 10,28 & 11,59 & 13,24 & 29,62 & 32,67 & 35,48 & 38,93 & 41,40\\
22 & 8,643 & 9,542 & 10,98 & 12,34 & 14,04 & 30,81 & 33,92 & 36,78 & 40,29 & 42,80\\
23 & 9,260 & 10,20 & 11,69 & 13,09 & 14,85 & 32,01 & 35,17 & 38,08 & 41,64 & 44,18\\
24 & 9,886 & 10,86 & 12,40 & 13,85 & 15,66 & 33,20 & 36,42 & 39,36 & 42,98 & 45,56\\
25 & 10,52 & 11,52 & 13,12 & 14,61 & 16,47 & 34,38 & 37,65 & 40,65 & 44,31 & 46,93\\
\addlinespace[2pt]
26 & 11,16 & 12,20 & 13,84 & 15,38 & 17,29 & 35,56 & 38,89 & 41,92 & 45,64 & 48,29\\
27 & 11,81 & 12,88 & 14,57 & 16,15 & 18,11 & 36,74 & 40,11 & 43,19 & 46,96 & 49,64\\
28 & 12,46 & 13,56 & 15,31 & 16,93 & 18,94 & 37,92 & 41,34 & 44,46 & 48,28 & 50,99\\
29 & 13,12 & 14,26 & 16,05 & 17,71 & 19,77 & 39,09 & 42,56 & 45,72 & 49,59 & 52,34\\
30 & 13,79 & 14,95 & 16,79 & 18,49 & 20,60 & 40,26 & 43,77 & 46,98 & 50,89 & 53,67\\
\addlinespace[2pt]
40 & 20,71 & 22,16 & 24,43 & 26,51 & 29,05 & 51,81 & 55,76 & 59,34 & 63,69 & 66,77\\
50 & 27,99 & 29,71 & 32,36 & 34,76 & 37,69 & 63,17 & 67,50 & 71,42 & 76,15 & 79,49\\
60 & 35,53 & 37,48 & 40,48 & 43,19 & 46,46 & 74,40 & 79,08 & 83,30 & 88,38 & 91,95\\
\bottomrule\end{tabular}
\end{center}
\end{document}
